% SHARD-ID: CA-12-LATTICE-BOOST-CURRENT-1D
% SHARD-TITLE: A One-Dimensional Boost Produces a Momentum Current
% SHARD-SUMMARY: Derives the nearest-neighbour identity that the first energy moment K=sum_j x_j h_j gives i[H,K] as a weighted sum of i[h_j,h_{j+1}].
% SHARD-SUMMARY: Interprets the uniform-spacing expression as the first lattice candidate for bulk momentum.
% SHARD-SUMMARY: Marks the derivation as local and conditional on the Hamiltonian-density decomposition and commutation range.
% SHARD-KEYWORDS: lattice boost, momentum current, Hamiltonian density, commutator, nearest neighbour, energy moment

\section{A One-Dimensional Boost Produces a Momentum Current}
\label{sec:lattice-boost-current-1d}

\subsection{Status, Sources, and Dependencies}

\textbf{Status: checked local derivation.}  The identity in this shard is not
quoted from a source.  It is derived from the local Hamiltonian-density ansatz
and checked by a small coefficient invariant.

Dependencies: CA-10, CA-11, and CONVENTIONS.md (h).

% Source: references/text/CFTFromLatticeFermions.txt:88 -- lattice Hamiltonian as sum of self-adjoint finite-support terms.
% Source: references/text/GaugingDefectsQuantumSpinSystems.txt:674 -- pairwise local interaction on neighbouring sites.
% Source: references/text/GaugingDefectsQuantumSpinSystems.txt:684 -- complete Hamiltonian given by summing local terms over sites.
% Checked: test/runtests.jl, "nearest-neighbour boost current coefficients".

\subsection{Setup}

Let a finite open one-dimensional lattice have sites or bonds labelled
\(j=1,\ldots,N\), real positions \(x_j\), and a self-adjoint local-density
decomposition
\[
  H=\sum_{j=1}^N h_j.
\]
Assume the density terms commute at distance greater than one:
\[
  [h_j,h_k]=0\qquad\text{if }|j-k|>1.
\]
This is the finite-range simplification appropriate for the first derivation;
longer-range densities only add longer-range edge terms.

Define the position-weighted energy moment
\[
  K_x=\sum_{j=1}^N x_j h_j.
\]
This is the lattice boost candidate in the sense of CA-10: a boost is treated as
a position-dependent time translation, and time translation is generated by
energy density.

\subsection{Commutator Calculation}

Compute
\[
\begin{aligned}
  i[H,K_x]
  &= i\left[\sum_j h_j,\sum_k x_k h_k\right]\\
  &= i\sum_{j,k}x_k[h_j,h_k].
\end{aligned}
\]
The terms with \(j=k\) vanish, and by the nearest-neighbour assumption only
pairs \((j,j+1)\) survive.  For each adjacent pair,
\[
\begin{aligned}
  i x_{j+1}[h_j,h_{j+1}]
  + i x_j[h_{j+1},h_j]
  &= i(x_{j+1}-x_j)[h_j,h_{j+1}].
\end{aligned}
\]
Hence
\[
  i[H,K_x]
  =
  \sum_{j=1}^{N-1}(x_{j+1}-x_j)\,i[h_j,h_{j+1}].
\]
Each \(i[h_j,h_{j+1}]\) is self-adjoint because \(h_j,h_{j+1}\) are
self-adjoint.

\subsection{Uniform Spacing}

If \(x_j=aj\), then \(x_{j+1}-x_j=a\) and
\[
  i[H,K_x]=a\sum_{j=1}^{N-1}i[h_j,h_{j+1}].
\]
For unit spacing this gives the candidate bulk momentum
\[
  P_{\mathrm{bulk}}^{\mathrm{lat}}
  :=
  \sum_{j=1}^{N-1}i[h_j,h_{j+1}],
\]
with endpoint terms deliberately visible.  Periodic chains require an extra
boundary convention because the position function \(j\) is not single-valued on
a circle.

\subsection{Interpretation}

The continuum bracket in CA-11 says that boost plus time translation gives a
spatial translation, up to sign convention.  The lattice calculation says that
the same operation applied to the first energy moment produces the commutator of
neighbouring energy densities.  Therefore \(i[h_j,h_{j+1}]\) is not guessed from
thin air: it is the local current forced by the first-moment ansatz.

\subsection{Boundary}

This is not yet a theorem that \(P_{\mathrm{bulk}}^{\mathrm{lat}}\) converges to
continuum momentum.  The missing ingredients are a scaling map, a domain for the
candidate operators, a state/sector sequence, and convergence of the commutator
relations.  CA-15 turns these into acceptance checks.
